\section{Ethics and responsible interpretation}\label{sec:ethics}

The survey of \cref{sec:design} asks people to respond to statements about Jews and about Israel, some of which many readers will find offensive and some of which express positions held in good faith by large parts of the public. Its respondents include Jewish people, Israelis, Palestinians, Muslims and people of every political view. Its results will be read in a debate in which the stakes are high and the parties distrust each other's evidence. This section sets out the safeguards the plan commits to and the limits on what the results can be used to show.

\subsection{Review, consent and debriefing}\label{sec:ethics-consent}

The protocol will be reviewed by an institutional review board before the pilot, and again before the main survey if the pilot leads to changes. The consent page states, before any question is shown, that the survey includes statements about Jews and about Israel, that some of these statements are offensive stereotypes, that respondents may skip any question or stop at any time without penalty, and that answers are confidential. Consent is recorded before the first item.

Two features of the design need explicit justification to the review board. First, respondents are shown stereotype statements. These are needed to measure the attitudes they express, and the item bank uses the smallest number that identification and precision allow (\cref{sec:design-items}). Second, the list experiment does not tell respondents which item on a list is of interest. This is indirect questioning, not deception: nothing false is stated, and the purpose of the technique is explained in the debriefing.

Every respondent who starts the survey, including those who stop early, sees a debriefing page. It states that the statements about Jews were included because they are common stereotypes whose prevalence the study measures, and that their inclusion does not mean the researchers regard them as true. It states that the Israel-related statements span views held across the political spectrum and that their inclusion implies no judgement about them. It explains the list experiment, gives contact details for the research team and the review board, and lists sources of support for respondents who found the content distressing. The pilot measures whether respondents report distress, and the debriefing is revised if it is inadequate.

\subsection{Data protection}\label{sec:ethics-data}

The survey data combine religion, political opinion, attitudes toward an ethnic and religious group, and geography. In the European Union, data revealing political opinions or religious beliefs are special categories of personal data \citep{gdpr2016}; any cross-national extension will be designed under that regime, and we treat the US data with the same care. The research team does not receive names, addresses or contact details; these stay with the panel provider, and respondents are linked across waves only through a pseudonymous identifier. Geography is held at the level needed for the small-area model and no finer. The analysis file is stored encrypted, with access limited to named analysts. Microdata are released only under restricted access, and openly only in synthetic form (\cref{sec:design-replication}). Released tables apply small-cell suppression and are checked for disclosure by differencing against other released tables \citep{hundepool2012}. Small-area results are released only as model-based estimates with intervals.

\subsection{Population associations are not individual diagnoses}\label{sec:ethics-individual}

The estimands of this paper describe populations. Four consequences follow, and every output will state them.
\begin{itemize}
  \item A respondent's estimated position on $G$ is a model-based summary with substantial uncertainty. It is used only to estimate population quantities. No individual score will be reported, and the methods are not a screening tool for individuals.
  \item An item's membership of the boundary set $B_\tau$ means that, \emph{in the surveyed population}, most of the common variance in responses to it is shared with responses to classic stereotypes. It does not mean that a particular person who endorses the item holds hostile attitudes toward Jews, or that a particular utterance with similar content is antisemitic in its context. Conversely, an item outside $B_\tau$ is not thereby shown to be acceptable in every context. Definitions of antisemitism classify speech and conduct in context, often with reference to intent (\cref{sec:bg-definitions}); our estimands do not.
  \item Differences in mean $G$ between ideological, religious, age or regional groups are differences between averages of overlapping distributions. They will always be reported with intervals and alongside the within-group spread, never as rankings of groups. They say nothing about any member of a group.
  \item Small-area estimates describe areas, not the people in them, and model-based estimates for small areas lean on the model. Estimates for named campuses will not be published (\cref{sec:design-sampling}), because the risk of stigmatising a student body outweighs the value of the estimate.
\end{itemize}
These points apply to every group, including groups whose members are themselves frequent targets of prejudice. The analysis plan does not single out any religious or ethnic group for comparison; religiosity enters as a measure of practice, not of affiliation (\cref{tab:design-groups}).

\subsection{Misuse in political debate}\label{sec:ethics-misuse}

The results bear on a contested question, and any of them can be quoted selectively. An item placed in $B_\tau$ can be cited as proof that a category of criticism of Israel is antisemitic; an item placed outside it can be cited as proof that it never is. Neither reading is supported. We commit to the following.
\begin{itemize}
  \item \emph{Publication whatever the results.} Every pre-registered estimand will be published, whichever position in the definitional debate it appears to support and including indeterminate classifications and null results (decision rule D5).
  \item \emph{No adjudication.} No output will describe the results as validating or refuting the IHRA working definition, the Jerusalem Declaration or the Nexus Document. Results will be presented item by item, with the definitional tags of \cref{sec:design-items}, so that readers can relate them to any definition themselves.
  \item \emph{Full context in every release.} Releases will state the population, the fieldwork dates, the threshold $\tau$ and the uncertainty of every classification. Plain-language summaries will say what the results do not show.
  \item \emph{Correction.} Errors found after publication will be corrected publicly, in the paper and in the replication package.
\end{itemize}

\subsection{Advisory input}\label{sec:ethics-advisory}

Before items are drafted, we will convene an advisory group whose members hold different views on the definitions of \cref{sec:bg-definitions}. It will include representatives of Jewish communal organisations with differing positions on those definitions; Palestinian, Arab and Muslim community organisations; scholars who have argued for different definitions; civil-liberties and academic-freedom specialists; and survey methodologists. The group reviews the item bank for balance and wording, assigns the definitional tags independently, reviews the consent and debriefing texts, and comments on how results are communicated. It has no veto over the results, and its members are not asked to endorse them. Membership, funding sources and any relationship between the research team and organisations with a stake in the results will be disclosed in every output.

\subsection{Neutrality}\label{sec:ethics-neutrality}

The Kornstein Lab exists to study antisemitism. Readers are entitled to ask whether that mission shapes the instrument or the analysis. We do not ask readers to rely on our intentions. The design builds in four protections. The threshold $\tau$, the decision rules and the item bank are fixed before data collection (\cref{sec:design}). The item bank covers every disputed category of Israel-related speech, and its balance is reviewed by advisers with opposing views (\cref{sec:ethics-advisory}). The estimands are symmetric: the same quantity is computed for every item, and an item can fall in or out of $B_\tau$ under the same rule. And the code, the synthetic data and the registered plan are public (\cref{sec:design-replication}), so that anyone can check that the analysis follows the plan. The paper's aim is to make an empirical sub-question of the boundary debate answerable by evidence that people on every side can inspect, not to settle the debate.
